\begin{table*}[t]
\centering
\small
\caption{G13 comparado aos modelos clássicos gerados com o mesmo $n$ e $m$ aproximado.}
\label{tab:modelos}
\begin{adjustbox}{max width=\linewidth}
\begin{tabular}{lrrrrrrrr}
\toprule
Modelo & n & m & grau máx & C médio & L & γ (bins log) & R² (bins log) & γ (MLE) \\
\midrule
G13 (maior componente) & 34.671 & 123.570 & 2.336 & 0,186 & 3,96 & 2,11 & 0,992 & 2,22 \\
Erdős–Rényi G(n,m) & 34.671 & 123.570 & 22 & $1,65 \times 10^{-4}$ & 5,52 & 1,46 & 0,195 & 2,92 \\
Barabási–Albert (m=4) & 34.671 & 138.668 & 620 & 0,00235 & 4,27 & 2,87 & 0,99 & 2,88 \\
Watts–Strogatz (k=8, p=0,1) & 34.671 & 138.684 & 12 & 0,473 & 8,07 & -5,03 & 0,361 & 5,71 \\
\bottomrule
\end{tabular}
\end{adjustbox}
\end{table*}
